% Code-coverage subsection.

\subsection{Code Coverage}
\label{sec:coverage}

% Coverage headline table.
% Single-column layout. Citations inline. Std rendered inline at
% \tiny (one step smaller than the table's \scriptsize cell font).
% Bolded mean marks column-best.

\begin{table}[!htbp]
\centering
\caption{Final line coverage mean $\pm$ std (\%) on the six
$(\mathrm{SUT, format})$ pairs. ``---'' indicates the tool is
not language-applicable.}
\label{tab:coverage-headline}
\setlength{\tabcolsep}{3pt}
\renewcommand{\arraystretch}{1.0}
\resizebox{\columnwidth}{!}{%
\begin{tabular}{l c c c c c c}
\toprule
& \multicolumn{3}{c}{\textbf{VCF}} & \multicolumn{3}{c}{\textbf{SAM}} \\
\cmidrule(lr){2-4} \cmidrule(lr){5-7}
Tool & htsjdk & vcfpy & noodles & htsjdk & biopy & seqan3 \\
\midrule
Pure-Random                              & $1.5$\,{\tiny$\pm0.1$}  & $2.7$\,{\tiny$\pm0.0$}  & $0.0$\,{\tiny$\pm0.0$}  & $2.2$\,{\tiny$\pm0.3$}  & $1.8$\,{\tiny$\pm1.2$}  & $78.3$\,{\tiny$\pm8.1$} \\
Jazzer~\citep{jazzer2024}                & $34.7$\,{\tiny$\pm0.0$} & ---                     & ---                     & $25.1$\,{\tiny$\pm0.2$} & ---                     & ---                     \\
Atheris~\citep{atheris2023}              & ---                     & $55.0$\,{\tiny$\pm0.0$} & ---                     & ---                     & $54.0$\,{\tiny$\pm0.5$} & ---                     \\
cargo-fuzz~\citep{cargofuzz2024}         & ---                     & ---                     & $22.9$\,{\tiny$\pm0.1$} & ---                     & ---                     & ---                     \\
libFuzzer~\citep{libfuzzer2024}          & ---                     & ---                     & ---                     & ---                     & ---                     & $\mathbf{98.5}$\,{\tiny$\pm0.6$} \\
AFL++~\citep{fioraldi2020aflpp}          & ---                     & ---                     & ---                     & ---                     & ---                     & $79.1$\,{\tiny$\pm0.0$} \\
\midrule
\textbf{\biotest{} (ours)}               & $\mathbf{45.7}$\,{\tiny$\pm1.7$} & $\mathbf{72.3}$\,{\tiny$\pm0.3$} & $\mathbf{37.1}$\,{\tiny$\pm0.9$} & $\mathbf{30.9}$\,{\tiny$\pm0.1$} & $\mathbf{54.3}$\,{\tiny$\pm0.0$} & $93.7$\,{\tiny$\pm0.3$} \\
\bottomrule
\end{tabular}}
\end{table}


\paragraph{Setup.}
We measure line coverage on six $(\mathrm{SUT, format})$ pairs,
scoped by the same target-class filter as the
mutation-adequacy metric of \S\ref{sec:mutation-adequacy}, so
coverage and mutation are graded on the same SUT surface. Line coverage (the fraction of source
lines in each parser executed at least once during a tool's
campaign) is collected by each language's standard coverage
instrumentation (specific backends are listed in
\S\ref{sec:implementation}). We report final-attainment line
coverage, the end-of-campaign value each tool reaches inside its
scope, because \biotest's Phase-D feedback loop
(Figure~\ref{fig:architecture}) terminates on its own stop-criteria
rather than at a fixed wall-clock tick. Within each row of
Table~\ref{tab:coverage-headline} the scope filter and
instrumentation backend are identical across columns. Line
percentages are directly comparable as final attainment.

\paragraph{Take-away.}
\biotest{} leads on five of six SUTs and trails libFuzzer by
$4.75$\,pp on seqan3/SAM. The split has a concrete structural
explanation. The five leading parsers expose declarative
validators (\texttt{INFO}, \texttt{FORMAT}, contig declarations
on VCF; CIGAR-shape, tag-type, and flag-bit checks on the SAM
parsers) that \biotest{}'s spec-grounded generator exercises
preferentially because each is named in the spec and retrieved
by LLM mining. Only seqan3/SAM retains the pattern Klees et al.\
and B\"ohme et al.~\cite{klees2018fuzz,boehme2020boosting}
report on byte-content-lenient parsers, where coverage-guided
byte mutators accumulate inputs reaching rare control-flow paths
more efficiently than spec-driven enumeration.
